% !TEX root = /home/eb/.vscode/extensions/qftrules.programme-colle/tmp/Exercice.tex

%%%%%%%%%%%%%%%%%%%%%%%%%%%%%%%%%%%
\begin{exo}[PC][2][TD]{Un exemple d'onde électromagnétique}
  \Notions{\item Onde électromagnétique dans le vide  \task Énergie électromagnétique}
	On étudie une onde électromagnétique dans le vide illimité, de longueur d'onde $\lambda=\SI{6e-7}{m}$, et dont le champ électrique est de la forme
	$$
	\udl{\E} =  \udl{E}_{x} \ex+\udl{E}_{y} \ey \quad \text { avec } \quad \udl{E}_{x}=E_{0} \exp \left[i\left(\omega t - \frac{K}{3}(2 x+2 y-z)\right)\right].
$$
\begin{quest}
	\item Identifier le domaine du spectre électromagnétique associé. Calculer la fréquence de l'onde à partir de la relation de dispersion. 
	\sol{
	$f = \frac{c}{l} = \SI{5e14}{Hz}$, lumière visible.
	}

	\item Exprimer le vecteur d'onde $\Vec{k}$ en fonction de la constante $K$, puis calculer la valeur numérique de $K$. 
  % Établir l'équation cartésienne d'un plan d'onde.
	\sol{
$\Vec{k} = \frac{K}{3}(2\ex + 2 \ey - \ez)$.
	}

	\item À partir de l'équation de Maxwell-Gauss, exprimer $\udl{E}_{y}$ en fonction de $\udl{E}_{x}$.

	\item Exprimer le champ magnétique $\udl{\B}$ de cette onde en fonction de $\udl{E}_{x}$.

	\item Exprimer la densité moyenne d'énergie électromagnétique $\moy{u_{\rm em}}$ et la moyenne temporelle du vecteur de Poynting $\moy{\Poy}$ de cette onde. Commenter.
\end{quest}
\end{exo}
%%%%%%%%%%%%%%%%%%%%%%%%%%%%%%%%%%%


%%%%%%%%%%%%%%%%%%%%%%%%%%%%%%%%%%%
\begin{exo}[PC][1][appli]{Polarisations}
  \Notions{\item Polarisation}
  Dans chaque cas, décrire entièrement la polarisation de l'onde électromagnétique dans le plan $z=0$.

  \hauteurLargeurCadreReponse		{8mm}{8cm}
  \pourcentageDeLaPartieAGauche{1}
  \initialisationPartieGauche
  
  % \begin{corrige}
    % 	On a $\vv{\ptO \ptM'}=r \left(\cos(\theta) \ex+\sin(\theta) \ey\right)$.
    % \end{corrige}
    \begin{enonce}
      $\E(t) = E_0 \cos(\omega t) \ex$
    \end{enonce}
        
    \begin{enonce}
      $\E(t) = E_0 \cos(\omega t) \ex - E_0 \sin(\omega t) \ey$
    \end{enonce}
    
    \begin{enonce}
      $\E(t) = \frac{E_0}{2}\cos(\omega t) \ex + \frac{E_0 \sqrt{3}}{2} \cos(\omega t) \ey$
    \end{enonce}

    \begin{enonce}
      $\E(t) = \frac{E_0}{2}\cos(\omega t) \ex + \frac{E_0 \sqrt{3}}{2} \sin(\omega t) \ey$
    \end{enonce}

  \initialisationPartieDroite
  \finalisationDuPartageDePage 
\end{exo}
%%%%%%%%%%%%%%%%%%%%%%%%%%%%%%%%%%%

%%%%%%%%%%%%%%%%%%%%%%%%%%%%%%%%%%%%%%%%%%%%%%%%%
\begin{exo}[PC][2][TD]{Superposition de deux OPPH polarisées rectilignement}
  \Notions{\item Onde électromagnétique dans le vide \task OPPH}
Dans le vide illimité, on considère deux OPPH de même pulsation $\omega$ et de même polarisation rectiligne suivant l'axe $(Oz)$, définies par leur champ électrique complexe
\begin{equation*}
\begin{cases}
\vv{\udl{E}_1}(M,t)=E_0\exp[i(\omega t-\vv{k_1}\cdot\vv{r})]\ez,\\
\vv{\udl{E}_2}(M,t)=E_0\exp[i(\omega t-\vv{k_2}\cdot\vv{r})]\ez,
\end{cases}
\end{equation*}
où les vecteurs d'onde vérifient $\vv{k_1} = k (\cos \alpha\, \ex + \sin \alpha \,\ey)$ et $\vv{k_2} = k (\cos \alpha \,ex - \sin \alpha \,\ey)$, avec $0 \leq \alpha \leq \pi$.
\begin{quest}
\item Déterminer les champs magnétiques complexes $\vv{\udl{B}_1}$ et $\vv{\udl{B}_2}$ associés.
\item Représenter les familles de vecteurs $(\Vec{k_1},\udl{\E}_1,\udl{\B}_1)$ et $(\Vec{k_2},\E*_2,\udl{\B}_2)$ dans l'espace.
\item Déterminer les champs $\E*$ et $\B*$ résultant de la superposition des deux ondes. De quelle nature est cette onde ? Quelle est sa structure ?
\end{quest}
\end{exo}
%%%%%%%%%%%%%%%%%%%%%%%%%%%%%%%%%%%%%%%%%%%%%%%%%

%%%%%%%%%%%%%%%%%%%%%%%%%%%%%%%%%%%%%%%%%%%%%%%%%
\begin{exo}[PC][1][colle]{Étude d'une OPPH}
On considère une OPPH qui se propage dans le vide illimité, caractérisée par son champ électrique complexe $\vv{\underline{E}}(M,t)=\udl{E}_x\,\vv{e_x}+\udl{E}_y\,\vv{e_y}$ où $\udl{E}_x=E_0\,\e^{i\varphi}$ avec $\textstyle\varphi=\frac{k}{\sqrt{6}}(x+2y+z)-\omega t$. La longueur d'onde vaut $\lambda=\SI{600}{nm}$.
\begin{quest}
\item Déterminer la fréquence, le vecteur d'onde et l'équation des plans d'onde.
\sol{
Fréquence $\nu=\frac{\omega}{2\pi}$, vecteur d'onde $\Vec{k} = \frac{k}{\sqrt{6}}(\ex + 2 \ey + \ez)$, plan d'onde de normale $\Vec{k}$.
}
\item Déterminer le champ électrique et donner sa polarisation.
\sol{
Que vaut $E_y$ ? On le cherche sous la forme d'une OPPH $E_y = \udl{E_{y,0}}e^{i\vphi}$. Maxwell-Gauss : $\dive \udl{\E} = \dr{E_x}{x} + \dr{E_y}{y} = E_0 \frac{k}{\sqrt{6}} e^{i\vphi} +  \udl{E_{y,0}} \frac{2k}{\sqrt{6}} e^{i\vphi} = 0 $ dans le vide, donc
$\udl{E_{y,0}} = -\frac{E_0}{2}.$
Ainsi,
$$
\udl{\E}(M,t) = E_0\,\e^{i\varphi}(\ex - \frac{1}{2}\ey) = E_0\,\e^{i\varphi}\Vec{\veps},
$$
où $\Vec{\veps}=\ex - \frac{1}{2}\ey$ est un vecteur constant : polarisation rectiligne.
}
\item Déterminer le champ magnétique.
\sol{
On détermine $\B$ à partir de $\E$ par Maxwell-Faraday. On a
$$
\rot \udl{\E} = \Vec{k} \pv \E = \frac{k}{\sqrt{6}}E_0\,\e^{i(\Vec{k} \ps \Vec{r} - \omega t)}(-\ez + \ey - \ex) = -\dr{\udl{\B}}{t},
$$
d'où
$$
\udl{\B}(M,t) = \frac{kE_0}{\sqrt{6}\omega}\,\e^{i(\Vec{k}\ps \Vec{r} - \omega t)} (-\ex + \ey - \ez).
$$
}
\item Faire une étude énergétique et commenter.
\sol{
Poynting $\Poy = \frac{\E\pv\B}{\mu_0}$ où $\E$ et $\B$ sont les champs \textbf{réels} !
}
\end{quest}
\end{exo}
%%%%%%%%%%%%%%%%%%%%%%%%%%%%%%%%%%%%%%%%%%%%%%%%%

%%%%%%%%%%%%%%%%%%%%%%%%%%%%%%%%%%%
\begin{exo}[PC][2][TD]{Onde électromagnétique cylindrique}
  \Notions{\item Onde électromagnétique dans le vide \task Analyse énergétique \task Coordonnées cylindriques}
  Une onde électromagnétique émise par des sources placées le long d'un axe $O z$ possède une structure cylindrique. En coordonnées cylindriques, le champ électrique d'une onde progressive divergente, monochromatique et polarisé suivant l'axe $O z$, possède une amplitude $E(r)$ fonction de la distance $r$ à l'axe, d'expression
  en notation complexe : 
  \begin{equation}
  \E*=E(r) \mathrm{e}^{i(\omega t-k r)} \ez,
  \end{equation}
  où la fonction $E(r)$ est supposé réelle.
  Il s'agit de déterminer la variation de $E(r)$ en fonction de $r$ sachant que la propagation s'effectue dans le vide où $\omega = kc$.
  \begin{quest}	
    \item Déterminer le champ magnétique $\B*$ à partir de l'équation de Maxwell-Faraday. Commenter.
    %  Commenter les deux termes qui le composent.
    \item Déterminer la valeur moyenne $\avg{\Poy}$ du vecteur de Poynting. En déduire la puissance moyenne $\mc{P}$ rayonnée par l'onde à travers un cylindre d'axe $(Oz)$ de hauteur $h$ et de rayon $r$.
    \item Que dire de cette puissance $\mc{P}$ ? En déduire $E(r)$ en fonction de $r$, de $\mc{P}$, et des constantes $\epz$ et $c$.
    \item Dans l'hypothèse $r \gg \lambda$ (champ lointain), exprimer les champs $\E*$ et $\B*$ et montrer que localement la structure de l'onde est celle d'une onde plane.
  \end{quest}
  
%   Il faut évaluer $\overrightarrow{\operatorname{rot}} \vec{E}$ pour l'identifier à $-\partial \vec{B} / \partial t$.
%   $$
%   \begin{aligned}
%   \overrightarrow{\operatorname{rot}} \underline{\vec{E}} & =\mathrm{e}^{i \omega t} \overrightarrow{\operatorname{rot}}\left(E(r) \mathrm{e}^{-i k r} \cdot \ez\right)=\mathrm{e}^{i \omega t}[E(r) \mathrm{e}^{-i k r} \underbrace{\overrightarrow{\operatorname{rot}} \ez}_{\overrightarrow{0}}+\overrightarrow{\operatorname{grad}}\left(E(r) \mathrm{e}^{-i k r}\right) \wedge \ez] \\
%   & =\mathrm{e}^{i \omega t}\left[E^{\prime}(r)-i k E(r)\right] \mathrm{e}^{-i k r} \cdot \vec{u}_r \wedge \ez=-\left(E^{\prime}(r)-i k E(r)\right) \mathrm{e}^{i(\omega t-k r)} \cdot \vec{u}_\theta
%   \end{aligned}
%   $$
  
%   $-\frac{\partial \vec{B}}{\partial t}=-i \omega \underline{\vec{B}}$ car $\vec{B}$ varie à la même pulsation.
%   d'où
  
%   $$
%   \underline{\vec{B}}=-\frac{1}{\omega}\left(i E^{\prime}(r)+k E\right) \mathrm{e}^{i(\omega t-k r)} \vec{u}_\theta
%   $$
  
%   $\underline{B}$ est polarisé rectilignement suivant $\vec{u}_\theta$ et est composé de deux termes, l'un en $-E(r) / c$ en phase avec $\underline{\vec{E}}$ et l'autre en $-i E^{\prime}(r) / \omega$ en quadrature avec $\underline{\vec{E}}$.
%   Cette structure respecte $\operatorname{div} \vec{B}=0\left(\vec{B}=B(r) \vec{u}_\theta\right)$, de même que celle de $\underline{\vec{E}}$ respecte $\operatorname{div} \vec{E}=0,\left(\vec{E}=E(r) \ez\right)$; l'onde est donc transverse puisque la direction de propagation en tout point est $\vec{u}_r$.
%   b) La valeur moyenne du vecteur de Poynting est donnée par
  
%   $$
%   <\vec{R}>=\frac{1}{2} \operatorname{Re}\left(\frac{\overrightarrow{\vec{E}} \wedge \vec{B}^*}{\mu_0}\right) \Longrightarrow
%   $$
  
  
%   $$
%   <\vec{R}>=\frac{1}{2} \varepsilon_0 c E^2(r) \vec{u}_r
%   $$
%   où il apparaît bien sûr que la composante de $\underline{\vec{B}}$ en quadrature avec $\underline{\vec{E}}$ ne contribule pas à la propagation de l'énergie : $\left(\underline{\vec{E}} \wedge \underline{\vec{B}}^*=E\left(\frac{E}{c}-\frac{i E^{\prime}}{\omega}\right)\right)$
% Le vecteur de Poynting est bien radial divergent.
% A travers un cylindre de rayon $r$ et de hauteur $h$, la puissance moyenne rayonnée est :

% $$
% P=\oiint<\vec{R}>\cdot \overrightarrow{\mathrm{d} S}=\frac{\varepsilon_0 c}{2} E^2(r) \cdot 2 \pi r h \Longrightarrow \quad P_h=\pi \varepsilon_0 c r E^2(r)
% $$

% c) Dans le vide considéré ici, il n'y a aucun phénomène dissipatif ; la puissance moyenne par unité de longueur est donc indépendante de $r$; elle correspond à la puissance d'émission
% d'où

% $$
% E(r)=\sqrt{\frac{P_h}{\pi \varepsilon_0 c r}}
% $$

% décroissance en $\frac{1}{\sqrt{r}}$ intermédiaire entre l'onde plane et l'onde sphérique.
% d) Avec $E=\frac{\alpha}{\sqrt{r}}$, il vient $\frac{E^{\prime}}{E}=-\frac{1}{2 r}$; l'amplitude du champ magnétique s'écrit alors $\quad \underline{B}(r)=-\frac{1}{\omega}\left(i E^{\prime}+k E\right)=-\left(1-\frac{i}{2 k r}\right) \frac{E}{c}$.
% Dans l'hypothèse du champ lointain, $r \gg \lambda \Longrightarrow k r \gg 1$ car $k=2 \pi / \lambda$, et donc $\underline{B}(r) \simeq-E / c$ devient réel (c'est donc le champ proche, lorsque $k r \ll 1$, qui est en quadrature avec $\underline{\vec{E}}$ ).

% $$
% \begin{aligned}
% & \text { d'où }\left\{\begin{array}{l}
% \underline{\vec{E}}=E(r) \mathrm{e}^{i(\omega t-k r)} \ez \\
% \underline{\vec{B}}=-\frac{E(r)}{c} \mathrm{e}^{i(\omega t-k r)} \vec{u}_\theta
% \end{array} \quad \text { avec } E(r)=\sqrt{\frac{P_h}{\pi \varepsilon_0 c r}}\right. \\
% & \text { d'où il apparaît que } \quad \vec{B}=\frac{\vec{u}_r \wedge \vec{E}}{c}
% \end{aligned}
% $$


% L'onde a donc localement la structure d'une onde plane progressive : $\left(\vec{u}_r, \vec{E}, \vec{B}\right)$ trièdre direct trirectangle, la restriction localement étant relative à la décroissance des champs en $1 / \sqrt{r}$.
\end{exo}
%%%%%%%%%%%%%%%%%%%%%%%%%%%%%%%%%%%


%%%%%%%%%%%%%%%%%%%%%%%%%%%%%%%%%%%%%%%%%%%%%%%%%%%%%%%%%%%%%%
\begin{exo}[PC][2][TD]{Pression de radiation sur une voile solaire}
  \Notions{\item Onde électromagnétique dans le vide \task  Modèle corpusculaire de la lumière}
Une voile solaire est un dispositif de propulsion permettant de se déplacer dans
l’espace à la manière d’un voilier. Les photons émis par le Soleil entrent en collision
avec la voile et lui cèdent leur quantité de mouvement, ce qui lui permet d’avancer.
%Compte tenu de la faible propulsion générée, le procédé ne permet pas de quitter la surface d’une planète (même dénuée d’atmosphère, et donc de friction). Il est en revanche utilisable sur un appareil ayant déjà atteint la vitesse de satellisation minimale, voire la vitesse de libération.
Plusieurs prototypes de petite taille ont
déjà été placés en orbite ou sont en cours de développement, comme par exemple
le démonstrateur Ikaros, lancé en
2010 par l’agence spatiale japonaise.

\begin{fig}[]
	\centering
	\includegraphics[width=0.4\textwidth]{ikaros.jpeg}
	\capt{Vue d'artiste de la voile solaire Ikaros.}
\end{fig}
On considère une OPPH qui se propage dans le vide illimité selon $\vv{e_x}$ caractérisée par son champ électrique polarisé rectilignement $\vv{E}(x,t)=E_0\cos(\omega t-kx)\vv{e_y}$.
\begin{quest}
\item Calculer l'intensité lumineuse $I$ définie comme la puissance moyenne qui traverse une surface unité perpendiculaire à la direction de propagation de l'onde. On exprimera $I$ en fonction de $\varepsilon_0$, $c$ et~$E_0$.
\sol{
relation de structure : $\B = \ex \pv \frac{\E}{c} = \frac{E_0}{c} \cos(\omega t - k x) \ez$. Le vecteur de Poynting vaut $\Poy = \frac{\E \pv \B}{\mu_0} = \frac{E_0^2}{c\mu_0} \cos(\omega t -k x)^2 \ex $. L'éclairement vaut ainsi
\eqbox{\mc{E} = \moy{\Pi}= \frac{E_0^2}{2c\mu_0} = \frac{1}{2} \epz c  E_0^2.}
}
\end{quest}

On assimile l'onde précédente à un faisceau de photons se propageant selon $\vv{e_x}$.
\begin{quest}
\item Soit $N_0$ le nombre de photons par unité de temps et de surface qui traverse un élément de surface perpendiculaire à $\vv{e_x}$. Exprimer $N_0$ en fonction de $\mc{E}$ et $\omega$.
% (et éventuellement de constantes dimensionnées).
\sol{
Chaque photon est l'énergie $E=h \nu = \hbar \omega$, donc $N_0 = \mc{E}/E = \frac{\epz c  E_0^2}{2 \hbar \omega} $.
}
\item L'onde arrive sur une surface plane perpendiculaire à $\vv{e_x}$, d'aire $S$, parfaitement absorbante. Par un bilan de quantité de mouvement sur la paroi lors de la collision avec un photon, déterminer la force $F$ subie par la paroi en fonction de $I$, $S$ et $c$.
\sol{
Chaque photon cède à la paroi une quantité de mouvement $p=\hbar k$ lors de l'absorption. La variation totale de quantité de mouvement par unité de temps vaut donc $N_0 p S = N_0 \hbar k S$. Cela est égal à la force $F$ subie par la paroi, soit
\eqbox{F = N_0 \hbar k S =  \frac{\epz c k E_0^2}{2 \omega} S =  \frac{1}{2}\epz E_0^2 S}
}
\item En déduire l'expression de la pression $P$ exercée sur la paroi en fonction de $\varepsilon_0$ et $E_0$, appelée pression de radiation. \AN~: calculer $P$ et $F$ pour $S=\SI{10}{m^2}$ et une intensité lumineuse solaire $I = \SI{1000}{W.m^{-2}}$.
\sol{
La pression de radiation $P = \frac{1}{2}\epz E_0^2 = u_e$ s'identifie à la densité volumique d'énergie. Pour un flux solaire $\Pi = \mc{E} = \SI{1000}{W.m^{-2}}$, on trouve $P = \frac{\Pi}{c} = \SI{3e-6}{Pa}$. Sur une voile de surface $S=\SI{10}{m^2}$, la force vaut $F = P S = \SI{3e-5}{N}$.
}
\end{quest}
\end{exo}
%%%%%%%%%%%%%%%%%%%%%%%%%%%%%%%%%%%%%%%%%%%%%%%%%%%%%%%%%%%%%%%%%%%%


%%%%%%%%%%%%%%%%%%%%%%%%%%%%%%%%%%%%%%%%%%%%%%%%%%%%%%%%%%%%%%
\begin{exo}[PSI][2]{Voile solaire}
Une voile solaire est un dispositif de propulsion permettant de se déplacer dans
l’espace à la manière d’un voilier. Les photons émis par le Soleil entrent en collision
avec la voile et lui cèdent leur quantité de mouvement, ce qui lui permet d’avancer.
%Compte tenu de la faible propulsion générée, le procédé ne permet pas de quitter la surface d’une planète (même dénuée d’atmosphère, et donc de friction). Il est en revanche utilisable sur un appareil ayant déjà atteint la vitesse de satellisation minimale, voire la vitesse de libération.
Plusieurs prototypes de petite taille ont
déjà été placés en orbite ou sont en cours de développement, comme par exemple
le démonstrateur Ikaros, lancé en
2010 par l’agence spatiale japonaise.

\begin{fig}[]
	\centering
	\includegraphics[width=0.4\textwidth]{ikaros.jpeg}
	\capt{Vue d'artiste de la voile solaire Ikaros.}
\end{fig}

On considère une voile solaire de surface $S$ modélisée par un conducteur parfait occupant le demi-espace $z>0$. Le rayonnement solaire est assimilé
à une onde plane progressive monochromatique (OPPM) de polarisation rectiligne. On suppose que la normale àsurface S est colinéaire à la direction de propagation de l’OPPM.


\begin{quest}

	\item Proposer une expression du champ électrique complexe de l’OPPM incidente sur la voile. Que vaut le champ électromagnétique dans un conducteur parfait ? En déduire l’existence d’une onde
réfléchie, dont on donnera l’expression du champ électrique. Que vaut alors le champ électrique dans le demi-espace $z<0$ ?

\sol{
$ \Vec{E i} = E 0 e^{i(\omega t - kz)} \Vec{e x}$ et $\Vec{E r} = -E 0 e^{i(\omega t + kz)} \Vec{e x}$ donc
$$
\Vec{E} = -2iE 0 \sin(kz) \Vec{e x}.
$$

}

	\item Calculer les champs magnétiques incident et réfléchis. En déduire la densité surfacique de courant $\Vec{j}_s$ parcourant la voile.

\sol{
$\Vec{B i} = \dfrac{E 0}{c}   e^{i(\omega t - kz)} \Vec{e y}$ et  $\Vec{B r} = \dfrac{E 0}{c}   e^{i(\omega t + kz)} \Vec{e y}$ donc
$$
\Vec{B} = \dfrac{2E 0}{c} \cos(kz) \Vec{e y}.
$$
La relation de passage donne
$$
\Vec{j_s} = \dfrac{2E 0}{\mu 0 c} e^{i\omega t} \ex.
$$
}

\item Quelle est l’expression de la force volumique magnétique ? Proposer une expression pour la force surfacique puis pour la force totale moyenne à laquelle est soumise la voile. Commenter
sa direction et son sens.

\sol{
$
\Vec{f_s} = \Vec{j_s} \wedge \Vec{B}(z=0)
$ donc
$$
\Vec{F} = 2 S \varepsilon_0 E_0^2 \Vec{e_z}.
$$
}

\end{quest}
$\begin{array}{lll}
\textsc{Données :}
& \text{relation de passage du champ électrique} & :  \Vec{E_2} - \Vec{E_1} = \frac{\sigma}{\varepsilon_0} \Vec{n} \\
& \text{relation de passage du champ magnétique} & :  \Vec{B_2} - \Vec{B_1} = \mu_0 \Vec{j}_s \wedge \Vec{n} \\
\end{array}
$

\end{exo}
%%%%%%%%%%%%%%%%%%%%%%%%%%%%%%%%%%%%%%%%%%%%%%%%%%%%%%%%%%%%%%%%%%%%
